\section{\'Etale morphisms. \'Etale coverings}
\label{I.4}

We shall admit everything we shall need concerning flat morphisms;
these facts will be proved later, if necessary\footnote{\Cf
Exp\ptbl\SourceRef{IV}{IV}.}.

\begin{definition}
\label{I.4.1}
\begin{enumerate}
\item[a)] Let $f\colon X\to Y$ be a morphism of finite type.
One says that $f$ is \emph{\'etale} at~$x$ if $f$ is flat at~$x$
and net at~$x$. One says that $f$ is \'etale if it is so at every
point. One says that $X$ is \'etale at~$x$ over~$Y$, or that it
is a $Y$-prescheme \'etale at~$x$, etc.
\item[b)] Let $f\colon A\to B$ be a local homomorphism. One says
that $f$ is \'etale, or that $B$ is \'etale over~$A$, if $B$ is
flat and unramified over~$A$.\footnote{\Cf the second thoughts in
III~\SourceRef{III.1.2}{1.2}.}
\end{enumerate}
\end{definition}

\begin{proposition}
\label{I.4.2}
For $B/A$ to be \'etale, it is necessary and sufficient that
$\widehat{B}/\widehat{A}$ be so.
\end{proposition}

In
% original p. 5
fact, this is true separately for ``net'' and for ``flat''.

\begin{corollary}
\label{I.4.3}
Let $f\colon X\to Y$ be of finite type, and $x\in X$. Whether
$f$ is \'etale at~$x$ depends only on the local homomorphism
$\cal{O}_{f(x)}\to\cal{O}_x$, and indeed only on the corresponding
homomorphism of completions.
\end{corollary}

\begin{corollary}
\label{I.4.4}
Suppose that the residue field extension $\kres(A)\to\kres(B)$
is trivial, or that $\kres(A)$ is algebraically closed. Then
$B/A$ is \'etale if and only if $\widehat{A}\to\widehat{B}$
is an isomorphism.
\end{corollary}

Combine flatness and~\ref{I.3.7}.

\begin{proposition}
\label{I.4.5}
Let $f\colon X\to Y$ be a morphism of finite type. Then the set
of points where it is \'etale is open.
\end{proposition}

Indeed, this is true separately for ``net'' and ``flat''.

This proposition shows that one can dispense with ``pointwise''
statements in studying morphisms of finite type which are \'etale
somewhere.

\begin{proposition}
\label{I.4.6}
\begin{enumerate}
\item[(i)] An open immersion is \'etale.
\item[(ii)] The composite of two \'etale morphisms is \'etale.
\item[(iii)] Extension of the base.
\end{enumerate}
\end{proposition}

Indeed, (i) is trivial, and for (ii) and (iii) it suffices to note
that this is true for ``net'' and for ``flat''. To tell the truth,
there are also corresponding statements for local homomorphisms
(without a finiteness condition), which in any case will have to
appear in the multiplodoque (beginning with the case: net).

\begin{corollary}
\label{I.4.7}
A cartesian product of two \'etale morphisms is likewise \'etale.
\end{corollary}

\begin{corollary}
\label{I.4.8}
Let $X$ and $X'$ be of finite type over~$Y$, and $g\colon X\to X'$
a $Y$-morphism. If $X'$ is unramified over~$Y$ and $X$ is \'etale
over~$Y$, then $g$ is \'etale.
\end{corollary}

Indeed, $g$ is the composite of the graph morphism
$\Gamma_g\colon X\to X\times_Y X'$, which is an open immersion
by~\ref{I.3.4}, and the projection morphism, which is \'etale
because it is obtained from the \'etale morphism $X\to Y$ by the
``change of base'' $X'\to Y$.

\begin{definition}
\label{I.4.9}
An \'etale (\resp net) covering of~$Y$ is a $Y$-scheme $X$ which
is finite over~$Y$ and \'etale (\resp net) over~$Y$.
\end{definition}

The first condition means that $X$ is defined by a coherent sheaf
of algebras $\cal{B}$ on~$Y$. The second then means that $\cal{B}$
is locally free on~$Y$ (\resp nothing at all), \emph{and} that,
in addition, for every $y\in Y$, the fiber
$\cal{B}(y)=\cal{B}_y\otimes_{\cal{O}_y}\kres(y)$ is a separable
algebra (= finite product of finite separable extensions)
over~$\kres(y)$.

\begin{proposition}
\label{I.4.10}
Let $X$ be a flat covering of~$Y$ of degree~$n$ (the definition
of this term deserved to appear in~\ref{I.4.9}), defined by a
coherent locally free sheaf $\cal{B}$ of algebras. One defines
in the well-known way the trace homomorphism $\cal{B}\to\cal{A}$
(which is a homomorphism of $\cal{A}$-modules, where
$\cal{A}=\cal{O}_Y$). For $X$ to be \'etale, it is necessary
and sufficient that the corresponding bilinear form
$\trace_{\cal{B}/\cal{A}}xy$ define an isomorphism of~$\cal{B}$
onto~$\cal{B}$, or, what amounts to the same, that the
\emph{discriminant section}
\[
d_{X/Y} = d_{\cal{B}/\cal{A}} \in
\Gamma(Y,\tbigwedge^n\check{\cal{B}}\otimes_{\cal{A}}
\tbigwedge^n\check{\cal{B}})
\]
be invertible, or finally that the discriminant ideal defined by
this section be the unit ideal.
\end{proposition}

Indeed, one is reduced to the case where $Y=\Spec(k)$, and then
this is a well-known separability criterion, trivial on passing
to the algebraic closure of~$k$.

\begin{remark}
We
% original p. 6
shall have a less trivial statement below, when one does not
assume a priori that $X$ is flat over~$Y$ but makes a normality
assumption.
\end{remark}

\section{The fundamental property of \'etale morphisms}
\label{I.5}

\begin{theorem}
\label{I.5.1}
Let $f\colon X\to Y$ be a morphism of finite type. For $f$ to be
an open immersion, it is necessary and sufficient that it be
an \emph{\'etale and radicial} morphism.
\end{theorem}

Recall that radicial means: injective, with purely inseparable
residue field extensions (and one may also recall that this means
that the morphism remains injective under every extension of the
base). Necessity is trivial; sufficiency remains. We shall give
two different proofs, the first shorter, the second more elementary.

1) A flat morphism is open, so one may suppose (replacing $Y$
by~$f(X)$) that $f$ is a \emph{homeomorphism} onto~$Y$.
Under any extension of the base, $f$ will remain flat, radicial,
and surjective, hence a homeomorphism, a fortiori closed.
Thus $f$ is \emph{proper}. Thus $f$ is \emph{finite}
(reference: Chevalley's theorem), defined by a coherent sheaf
$\mathcal{B}$ of algebras. $\mathcal{B}$ is locally free;
moreover, by the hypothesis it has rank~$1$ everywhere, so
$X=Y$, qed.

2) One may suppose that $Y$ and $X$ are \emph{affine}. One is
further easily reduced to proving the following: if $Y=\Spec(A)$,
with $A$ local, and if $f^{-1}(y)$ is nonempty ($y$ being the
closed point of~$Y$), then $X=Y$ (indeed, this will imply that
every $y\in f(X)$ has an open neighborhood $U$ such that $X|U=U$).
We shall have $X=\Spec(B)$; we want to prove that $A=B$. But for
this one is reduced to proving the analogous assertion on
replacing $A$ by~$\hat A$, and $B$ by $B\otimes_A\hat A$
(taking into account that $\hat A$ is faithfully flat over~$A$).
One may therefore suppose that $A$ is \emph{complete}. Let $x$
be the point above~$y$; by Corollary~\ref{I.2.2}, $\mathcal{O}_x$
is finite over~$A$ and hence (being flat and radicial over~$A$)
is identical to~$A$. Thus $X=Y\amalg X^\prime$ (disjoint sum).
Since $X$ is radicial over~$Y$, $X^\prime$ is empty. We are done.

\begin{corollary}
\label{I.5.2}
Let $f\colon X\to Y$ be a morphism which is a \emph{closed
immersion} and is \emph{\'etale}. If $X$ is connected, $f$ is
an isomorphism of~$X$ onto a connected component of~$Y$.
\end{corollary}

Indeed, $f$ is also an open immersion. One deduces:

\begin{corollary}
\label{I.5.3}
Let
% original p. 7
$X$ be a net $Y$-scheme, with $Y$ connected. Then every section
of~$X$ over~$Y$ is an isomorphism of~$Y$ onto a connected
component of~$X$. There is therefore a one-to-one correspondence
between the set of these sections and the set of connected
components $X_i$ of~$X$ such that the projection $X_i\to Y$
is an isomorphism (or, equivalently by~\ref{I.5.1}, is surjective
and radicial). In particular, a section is known when its value
at one point is known.
\end{corollary}

Only the first assertion requires a proof; by~\ref{I.5.2}, it
suffices to observe that a section is a closed immersion
(since $X$ is separated over~$Y$) and is \'etale by~\ref{I.4.8}.

\begin{corollary}
\label{I.5.4}
Let $X$ and $Y$ be two preschemes over~$S$, with $X$ net and
separated over~$S$ and $Y$ connected. Let $f$, $g$ be two
$S$-morphisms from~$Y$ to~$X$, and $y$ a point of~$Y$; suppose
that $f(y)=g(y)=x$ and that the residue field homomorphisms
$\kres(x)\to\kres(y)$ defined by $f$ and~$g$ are identical
(``$f$ and~$g$ agree geometrically at~$y$''). Then $f$ and~$g$
are identical.
\end{corollary}

This follows from~\ref{I.5.3} by reducing to the case where
$Y=S$, replacing $X$ by $X\times_S Y$.

Here is a particularly important variant of~\ref{I.5.3}.

\begin{theorem}
\label{I.5.5}
Let $S$ be a prescheme, $X$ and $Y$ two $S$-preschemes, $S_0$
a closed subprescheme of~$S$ with the same underlying space
as~$S$, and $X_0=X\times_S S_0$ and $Y_0=Y\times_S S_0$ the
``restrictions'' of~$X$ and~$Y$ to~$S_0$. Suppose that $X$ is
\'etale over~$S$. Then the natural map
\[
\Hom_S(Y,X)\to\Hom_{S_0}(Y_0,X_0)
\]
is \emph{bijective}.
\end{theorem}

One is again reduced to the case where $Y=S$, and then this
follows from the ``topological'' description of the sections
of~$X/Y$ given in~\ref{I.5.3}.

\begin{scholium}
This result contains an assertion of \emph{uniqueness} and
\emph{existence} of \emph{morphisms}. It can also be expressed
(when $X$ and~$Y$ are both taken \'etale over~$S$) by saying that
the functor $X\mto X_0$ from the category of \'etale $S$-schemes
to the category of \'etale $S_0$-schemes is \emph{fully faithful},
\ie establishes an \emph{equivalence} of the first with a
\emph{full subcategory} of the second. We shall see below that
it is even an equivalence of the first and the second
(which will be a theorem on the \emph{existence of \'etale
$S$-schemes}).
\end{scholium}

The
% original p. 8
following form of~\ref{I.5.5}, more general in appearance, is
often convenient:

\begin{corollary}[``Theorem on extending liftings'']
\label{I.5.6}
Consider a commutative diagram
\[
\xymatrix{
X\ar[d]&\ar[l]Y_0\ar[d]\\
S&\ar[l]Y
}
\]
of morphisms, where $X\to S$ is \'etale and $Y_0\to Y$ is
a bijective closed immersion. Then one can find a unique
morphism $Y\to X$ which makes the two corresponding triangles
commutative.
\end{corollary}

Indeed, replacing $S$ by~$Y$ and $X$ by $X\times_S Y$, one is
reduced to the case where $Y=S$, and then this is the special
case of~\ref{I.5.5} for $Y=S$.

Let us also point out the following immediate consequence
of~\ref{I.5.1} (which we did not give as Corollary~1 so as not
to interrupt the line of ideas developed after~\ref{I.5.1}):

\begin{proposition}
\label{I.5.7}
Let $X$, $X^\prime$ be two preschemes of finite type and flat
over~$Y$, and let $g\colon X\to X^\prime$ be a $Y$-morphism.
For $g$ to be an open immersion (\resp an isomorphism), it is
necessary and sufficient that, for every $y\in Y$, the induced
morphism on the fibers
\[
g\otimes_Y\kres(y)\colon X\otimes_Y\kres(y)\to X^\prime
\otimes_Y\kres(y)
\]
be so.
\end{proposition}

It suffices to prove sufficiency; since the assertion holds for
surjectivity, one is reduced to the case of an open immersion.
By~\ref{I.5.1}, one must check that $g$ is \emph{radicial},
which is trivial, and that it is \emph{\'etale}, which follows
from Corollary~\ref{I.5.9} below.

\begin{corollary}
\label{I.5.8}
(Should be moved to \No\ref{I.3}.) Let $X$ and $X^\prime$ be
two $Y$-preschemes, $g\colon X\to X^\prime$ a $Y$-morphism,
$x$ a point of~$X$, and $y$ its projection on~$Y$. For $g$ to
be quasi-finite (\resp net) at~$x$, it is necessary and sufficient
that $g\otimes_Y\kres(y)$ be so.
\end{corollary}

Indeed, the two algebras over~$k\big(g(x)\big)$ which one must
examine to check that one indeed has a morphism quasi-finite
\resp net at~$x$ are the same for~$g$ and $g\otimes_Y\kres(y)$.

\begin{corollary}
\label{I.5.9}
With
% original p. 9
the notation of~\ref{I.5.8}, suppose that $X$ and~$X^\prime$
are flat and of finite type over~$Y$. For $g$ to be flat
(\resp \'etale) at~$x$, it is necessary and sufficient that
$g\otimes_Y\kres(y)$ be so.
\end{corollary}

For ``flat'', the statement is included only as a reminder;
it is one of the fundamental flatness criteria\footnote{\Cf
IV~\SourceRef{IV.5.9}{5.9}.}. For \'etale, it follows from this,
taking~\ref{I.5.8} into account.

\section{Application to \'etale extensions of complete local rings}
\label{I.6}

This no.\ is a special case of results on formal preschemes
which will have to appear in the multiplodoque. Nevertheless,
one gets by here with less effort, \ie without the explicit
local determination of \'etale morphisms in
\No\ref{I.7} (using the Main Theorem). This is perhaps
sufficient reason to keep the present no.\ (even in the
multiplodoque) in this place.

\begin{theorem}
\label{I.6.1}
Let $A$ be a complete local ring (noetherian, of course), with
residue field~$k$. For every $A$-algebra~$B$, let
$R(B)=B\otimes_A k$, regarded as a $k$-algebra; it therefore
depends functorially on~$B$. Then $R$ defines an
\emph{equivalence} of the category of $A$-algebras
\emph{finite and \'etale over~$A$} with the category of
\emph{separable algebras of finite rank} over~$k$.
\end{theorem}

First of all, the functor in question is fully faithful, as
follows from the more general fact:

\begin{corollary}
\label{I.6.2}
Let $B$, $B^\prime$ be two $A$-algebras finite over~$A$. If $B$
is \'etale over~$A$, then the canonical map
\[
\Hom_{\textup{$A$-alg}}(B,B^\prime)\to
\Hom_{\textup{$k$-alg}}\big(R(B),R(B^\prime)\big)
\]
is bijective.
\end{corollary}

One is reduced to the case where $A$ is artinian (by replacing
$A$ by~$A/\goth{m}^n$), and then this is a special case
of~\ref{I.5.5}.

It remains to prove that for every finite separable $k$-algebra
(why not say: \'etale, it is shorter) $L$, there exists a $B$
\'etale over~$A$ such that $R(B)$ is isomorphic to~$L$. One may
suppose that $L$ is a separable extension of~$k$; as such, it
has a generator~$x$, \ie is isomorphic to an algebra $k[t]/Fk[t]$
where $F\in k[t]$ is a monic polynomial. One lifts~$F$ to a
monic polynomial $F_1$ in $A[t]$, and takes $B=A[t]/F_1A[t]$.
